\chapter{Introduction}

This document aims to provide a theoretical explanation for the finite element method (FEM) for solid mechanics problems. The document begins by presenting all the necessary theoretical background, including the definition of many kinematics and stress related variables. From those, the available material constitutive models are presented, which are used to relate the kinematics and stress variables. Then, equation to be solved is derived by using the principle of virtual work. It is then discretized by using the finite element method. Finally, both solution methods are explained. Addionally, the document will present an overview of the code structure and all of its components.

This project is based on a academic projet I made in 2025 at Centrale Nantes. Based on that, the approach has been adapted to be used in a more general way, and the code has been rewritten for better comprehension, modularity and optimize performance - obtained mainly through the use of numpy array instead of loops. Addionally, the book \textit{"Nonlinear Continuum Mechanics for Finite Element Analysis"} by Javier Bonet and Richard D. Wood was used for theoretical reference.
